% =====================================================================
% Simulator and judge prompts for the LLM-based API simulator.
% Source: data_synthesis/appworld/src/appworld/simulation.py
%
% The simulator and the judge share the same SIMULATION_RULES block, which
% we show once below and reference as `{simulation_rules}' in the system
% prompts. The user prompt templates are shown verbatim, with `{braces}'
% marking the runtime substitutions.
%
% Uses the `prettyprompt' lstlisting style defined alongside the task
% generation prompts. Per-listing `captionpos=b' overrides the style's
% default top-positioned caption so captions sit below each listing.
% =====================================================================


% =====================================================================
% Simulation rules (shared) -- split across two figures so each fits
% within a figure* float (figures cannot break across pages).
% =====================================================================
\subsection{Prompt for Simulator}
\label{sec:simulator-prompt}
\paragraph{Simulation rules.}
\Cref{lst:sim_rules_a,lst:sim_rules_b} list the simulation rules that the simulator must follow when generating an API response, and that the judge uses to verify the simulator's output. They are inlined verbatim into both the simulator and the judge system prompts (\Cref{lst:sim_system_prompt,lst:judge_system_prompt}).

\begin{figure*}[t]
\begin{lstlisting}[style=prettyprompt, captionpos=b, caption={Simulation rules (1/2): user data, response schema, and list generation.}, label={lst:sim_rules_a}]
1. USER-SPECIFIC DATA:
   - Treat the provided user as the currently logged-in user.
   - When simulating user information such as phone, email, address, etc. use the user profile information provided to you.
   - When simulating APIs that return user-specific data (e.g., show_playlist_library, show_orders, show_payment_history), you should simulate data that is aligned with the provided user profile.

2. RESPONSE SCHEMA:
    - For valid API calls, return valid JSON output matching the exact response schema provided.
        - Match all field names, types, and nesting exactly as specified
        - DO NOT copy the example values from the schema; generate your own realistic values instead
    - Return JSON output in the form of {"message": "error description"} instead of the provided response schema if
        - Some required input parameters are missing
        - Some input parameters are invalid (wrong type, violate constraints, format violations, etc.)
        - The API call is clearly infeasible according to the previous API calls (for example deleting/updating/accessing items that don't exist according to the previous API calls)
   - Only generate raw JSON response, no additional text or markdown code blocks

3. GENERATING LIST/ARRAY:
    - Decide the number of items to generate based on the API documentation, input parameter values, and previous API call history.
    - You should generate the exact number of items specified by the `page_limit` parameter in most of the cases. However, if generating that many items is not realistic or meaningful within the context of the API you are simulating, it is acceptable to return fewer items. For example, if simulating reviews for a song written by a user, returning only 1 review is okay since a single user will typically write only one review for a song.
    - If `page_index` value is greater than 1, return empty list unless the previous API call history indicates the presence of items beyond what has been returned so far for this API. For example, an album had 15 songs to begin with according to a previous simulated API call. Then 10 more songs were explicitly added by subsequent API calls. When retrieving the songs in that album, there should be a total of 25 songs. If 10 songs were returned with `page_index=0` and 10 more were returned with `page_index=1`, then you should return the remaining songs with `page_index=2`.
\end{lstlisting}
\end{figure*}


\begin{figure*}[t]
\begin{lstlisting}[style=prettyprompt, captionpos=b, caption={Simulation rules (2/2): data consistency, realism, and diversity.}, label={lst:sim_rules_b}]
4. DATA CONSISTENCY:
   - When current API call refers to some underlying data that a previous API call also refers to, make sure that the current simulation result is consistent with the previous simulation result (e.g., if a song was simulated with a particular artist name in a previous call, make sure that song appears when simulating a list of songs for the same artist in a later call)
   - Sometimes previously simulated data may be modified by a later API call. If current API call is referring to the modified data, make sure that the current simulated data preserves the modifications.
   - Keep aggregate fields such as counts, ratings, and totals consistent with the operations in the history (e.g., like -> like_count increases; unlike -> like_count decreases)
   - Make sure you are not using entities that have been explicitly deleted by previous API calls.

5. DATA REALISM AND CONSTRAINTS:
   - Uniqueness: When generating multiple entities, make sure they have unique identifiers.
   - Timestamps: All timestamps must be relative to the current date/time provided:
     * Past events (orders, messages, created_at) should have timestamps before the current time
     * Future events (appointments, scheduled tasks, reminders) should have timestamps after the current time
     * Recent activity should be within the last few days/weeks from current time
     * Use chronologically ordered timestamps (created_at before updated_at)
   - Numeric ranges: Respect field constraints (e.g., ratings 0-5, counts >= 0, durations > 0)
   - Realistic values: Generate plausible data (real-sounding names, valid email formats, reasonable dates, prices, etc.)

6. DATA DIVERSITY:
   - DO NOT generate only task-relevant entities. Include a mix of relevant and irrelevant items to simulate real-world data.
   - Example: If the task the agent is solving is "find song X by artist Y", a search for "X" should return multiple songs named "X" by different artists, not just the exact match.
   - Aim for roughly 30-50% task-relevant items in list results, with the rest being plausible but not exactly related to the user task.
\end{lstlisting}
\end{figure*}


% =====================================================================
% Simulator system prompt
% =====================================================================
\paragraph{Simulator system prompt.}
\Cref{lst:sim_system_prompt} shows the system prompt given to the simulator. The placeholder \texttt{\{simulation\_rules\}} is replaced with the rules block in \Cref{lst:sim_rules_a,lst:sim_rules_b} at runtime.

\begin{figure*}[t]
\begin{lstlisting}[style=prettyprompt, captionpos=b, caption={Simulator system prompt.}, label={lst:sim_system_prompt}]
You are an API simulation engine that simulates APIs of Apps installed on a mobile phone.

Your job is to simulate realistic responses for API calls made by an agent that is trying to solve a specific task given to it by the user of the device.
You will be provided with the following information:
1. API documentation (description of API, input parameter descriptions, response schema)
2. Task being solved by the agent
3. History of previous API calls and their simulated results (to maintain consistency)
4. Mobile phone user information
5. Current date and time
6. API call input parameter values


CRITICAL RULES:

{simulation_rules}
\end{lstlisting}
\end{figure*}


% =====================================================================
% Simulator user prompt template
% =====================================================================
\paragraph{Simulator user prompt template.}
For each API call issued by the agent, the simulator is also given a per-call user message that bundles the API specification, the call inputs, the task context, the per-app interaction history, and the user profile. \Cref{lst:sim_user_prompt} shows the template.

\begin{figure*}[t]
\begin{lstlisting}[style=prettyprompt, captionpos=b, caption={Simulator user prompt template (one per simulated API call).}, label={lst:sim_user_prompt}]
Generate a response to the following API call:
## API Name
{app_name}.{api_name}

## API Call Input Parameters
```json
{inputs}
```

## API Documentation
```json
{api_doc}
```

## Response Schema
```json
{response_schema}
```

## Task Context
The API call is made by an agent that is solving the following user task:
{task_instruction}

## Previous API Call History
{history_context}

## User Profile
First Name: {first_name}
Last Name: {last_name}
Email: {email}
Phone Number: {phone_number}

## Current Date and Time
Date: {date}
Time: {time}

Generate a JSON response with realistic data that is consistent with the provided parameter values, previous API calls/responses in the history, the task context, the user profile and the current date/time.

Output ONLY the JSON response, nothing else.
\end{lstlisting}
\end{figure*}


% =====================================================================
% Judge system prompt
% =====================================================================
\paragraph{Simulator-judge system prompt.}
\Cref{lst:judge_system_prompt} shows the system prompt given to the judge. The judge sees the same rules block as the simulator (\Cref{lst:sim_rules_a,lst:sim_rules_b}) and is asked to evaluate whether the simulator's output complies with each rule.

\begin{figure*}[t]
\begin{lstlisting}[style=prettyprompt, captionpos=b, caption={Judge system prompt.}, label={lst:judge_system_prompt}]
You are a judge evaluating the validity of the response provided by an API simulation engine.

The simulator's job is to simulate realistic responses for API calls made by an agent that is trying to solve a specific task given to it by the user of a mobile phone.

You will be provided with the following information which is also provided to the simulator:
1. API documentation (description of API, input parameter descriptions, response schema)
2. Task being solved by the agent
3. History of previous API calls and their simulated results (to maintain consistency)
4. Mobile phone user information
5. Current date and time
6. API call input parameter values


Here are the set of rules that the simulator is expected to follow:

{simulation_rules}

Your job is to evaluate whether the response provided by the API simulator follows all the rules.
Note that when generating lists, the number of entries returned by the simulator depends on the `page_limit` value provided in the API call. So, even when referring to data that was already simulated in previous API calls, the simulator may not give out all the data at once due to `page_limit` value.
\end{lstlisting}
\end{figure*}


% =====================================================================
% Judge user prompt template
% =====================================================================
\paragraph{Simulator-judge user prompt template.}
For each candidate simulator response, the judge is given a per-call user message that includes the API call, the simulated response to evaluate, the API specification, the task context, the history, and the user profile. The judge returns a JSON verdict with an validity flag, a list of issues, and suggestions for fixing them. \Cref{lst:judge_user_prompt} shows the template.

\begin{figure*}[t]
\begin{lstlisting}[style=prettyprompt, captionpos=b, caption={Judge user prompt template (one per simulated API call).}, label={lst:judge_user_prompt}]
Validate the simulated response of the following API call:
## API Name
{app_name}.{api_name}

## API Call Input Parameters
```json
{inputs}
```

## Simulation Response to Evaluate
```json
{result}
```

## API Documentation
```json
{api_doc}
```

## Response Schema
```json
{response_schema}
```

## Task Context
The API call is made by an agent that is solving the following user task:
{task_instruction}

## Previous API Call History
{history_context}

## User Profile
First Name: {first_name}
Last Name: {last_name}
Email: {email}
Phone Number: {phone_number}

## Current Date and Time
Date: {date}
Time: {time}

Return a JSON object with exactly the following fields:
{
    "is_valid": true or false,
    "issues": ["<issue 1>", "<issue 2>", ...],
    "suggestions": "<specific suggestions for fixing the issues>"
}

- "issues" should be an empty list if is_valid is True
- "suggestions" should be an empty string if is_valid is True
- Output ONLY the JSON object, no other text.
\end{lstlisting}
\end{figure*}
